\documentclass[11pt]{article}

\usepackage[T1]{fontenc}
\usepackage{lmodern}
\usepackage[a4paper,margin=28mm]{geometry}
\usepackage{amsmath,amssymb,amsthm,mathtools}
\usepackage{microtype}
\usepackage{enumitem}
\usepackage{xcolor}
\usepackage[colorlinks=true,linkcolor=blue!45!black,citecolor=blue!45!black,
  urlcolor=blue!55!black]{hyperref}
\setlength{\emergencystretch}{2em}

\newtheorem{theorem}{Theorem}[section]
\newtheorem{proposition}[theorem]{Proposition}
\newtheorem{lemma}[theorem]{Lemma}
\newtheorem{corollary}[theorem]{Corollary}
\theoremstyle{definition}
\newtheorem{definition}[theorem]{Definition}
\theoremstyle{remark}
\newtheorem{remark}[theorem]{Remark}

\newcommand{\C}{\mathbb C}
\newcommand{\R}{\mathbb R}
\newcommand{\N}{\mathbb N}
\newcommand{\cA}{\mathcal A}
\newcommand{\cD}{\mathcal D}
\newcommand{\cF}{\mathcal F}
\newcommand{\cG}{\mathcal G}
\newcommand{\cH}{\mathcal H}
\newcommand{\cK}{\mathcal K}
\newcommand{\cP}{\mathcal P}
\newcommand{\cQ}{\mathcal Q}
\newcommand{\cV}{\mathcal V}
\newcommand{\cZ}{\mathcal Z}
\newcommand{\diff}{\partial}
\newcommand{\Res}{\operatorname*{Res}}
\newcommand{\Graph}{\operatorname{Graph}}

\title{The Arithmetic Julia Normal-Cone Law:\\
From the Prime Logarithmic Clock to the Complete Off-Critical Riemann Divisor}
\author{Michael Haufschild}
\date{}

\begin{document}

\maketitle

\begin{abstract}
We construct a source-defined singular completion of the prime logarithmic clock and prove that
one normal coordinate has two exact arithmetic boundary faces.  Let
\(k_p=(1,p^{-1},p^{-2},\ldots)\in\ell^2(\N_0)\).  At every finite trace order \(m\), the graph
generated by the actual-prime packets
\[
 \sum_p d_pk_p,\qquad
 \sum_p d_p(\log p)p^{-m}k_p
\]
is dense in the full product Hilbert space.  Every bounded realization which recovers the second
coordinate packetwise therefore has operator norm tending to infinity.

We then sew the prime aperture \(q=p^{-x}\) to the archimedean characteristic
\(\phi_\infty(t)=(\cosh t)^{-1/2}\).  The resulting scalar \(c_x(p)\) is a strict contraction and
has the exact fourth-power law
\[
 c_x(p)^4=\frac{(1+q)^4}{8(1+q^4)}.
\]
Its positive Julia defect \(d_x(p)=\sqrt{1-c_x(p)^2}\) is the unique minimal
orientation-preserving two-channel completion of the scalar source cell.  The normalized normal
coordinate
\[
 \ell_x(p)=\frac{2}{\sqrt3}\frac{d_x(p)}x
\]
converges to the literal logarithmic clock \(\log p\).  On every finite prime packet this Julia
coordinate converges to the dense-graph clock above, so the cellwise construction and the global
Hardy obstruction are one singular completion rather than unrelated identities.

Squaring the same normal coordinate against a canonical soft prime-power source produces the
Mellin orbit
\[
 M_u(n)=\frac{\Lambda(n)}n e^{-u\log n}.
\]
The Euler face \(u=-1\) is the von Mangoldt Dirichlet carrier
\(-\zeta'/\zeta\), which we prove hypertranscendental over the differential field
\[
 \C(s)\langle P(rs):r\ge1\rangle_{\rm diff}
\]
generated by the coordinate and every prime-zeta valuation layer simultaneously.  The critical
face \(u=-\tfrac12\) is \(\Lambda(n)/\sqrt n\).  A fixed-width two-edge explicit-formula
scalarization of this face, completed only by the canonical polar and Gamma terms, yields a
single coefficient sequence whose Dirichlet continuation has simple poles exactly at the
shifted strict-right zeros of the completed zeta function; normalized residues recover analytic
multiplicity.  Thus the source determines the complete off-critical Riemann divisor, and its
continuation is holomorphic on \(\Re s>1\) exactly under the Riemann hypothesis.

Finally, the recovered divisor has a canonical infinitesimal defect index.  Every strict-right
upper Schwarz pair contributes exactly \(8\pi\), independently of height and transverse
displacement.  The total nonnegative index vanishes exactly under the Riemann hypothesis.  The
result does not prove that it vanishes.  It opens a lossless route from actual-prime Hilbert
geometry to the complete divisor while proving that bounded Hardy completions and finite
differential algebra cannot carry the required information.
\end{abstract}

\medskip
\noindent\textbf{Keywords.}
Riemann hypothesis; Julia operator; prime logarithmic clock; hypertranscendence;
von Mangoldt function; Weil explicit formula; divisor tomography; defect index.

\medskip
\noindent\textbf{2020 Mathematics Subject Classification.}
Primary 11M26, 11J91; Secondary 47A20, 12H05, 11M06.

\section{Introduction}\label{sec:intro}

Write
\[
 \Xi(s)=\frac12s(s-1)\pi^{-s/2}\Gamma(s/2)\zeta(s).
\]
The explicit formula gives a precise dictionary between primes and the zero divisor of
\(\Xi\), while Weil positivity turns the Riemann hypothesis into a sign condition on a large
test-function space \cite{Weil1952,Guinand1949,Barner1981,Bombieri2000,Bombieri2003}.  Spectral,
semilocal, and Hilbert-space programs seek a smaller object which retains enough of that
dictionary to force the zeros onto the critical line
\cite{Connes1999,ConnesConsani2021,ConnesConsani2023,CCM2025,Suzuki2025,Suzuki2026}.
The unresolved issue is not merely how to encode a zero condition.  It is how to complete the
prime source without discarding the logarithmic direction which differentiation and the
explicit formula require.

This paper gives an exact answer to the completion problem.  The answer is necessarily singular.
In the Hardy pivot \(\cH=\ell^2(\N_0;\R)\), let
\[
 k_p=(1,p^{-1},p^{-2},\ldots)
\]
for each prime \(p\), and for finitely supported real packets \(d=(d_p)\) put
\[
 C(d)=\sum_p d_pk_p,\qquad
 L_m(d)=\sum_p d_p(\log p)p^{-m}k_p.                 \tag{1.1}
\]
For every \(m\ge0\),
\[
 \overline{\Graph(C,L_m)}=\cH\oplus\cH.              \tag{1.2}
\]
Thus the source and its logarithmic clock become maximally decoupled after closure.  No
nonzero bounded postprocessing repairs this.  More strongly, if bounded operators recover
\(L_m\) on every finite packet, their norms tend to infinity.

Dense graph saturation identifies what a successful completion must do, but it does not
construct one.  The construction comes from a different geometry.  For aperture \(x>0\) and
\(q=p^{-x}\), combine the M\"obius--Szeg\H{o} prime cell with the Gamma characteristic
\(\phi_\infty(t)=(\cosh t)^{-1/2}\):
\[
 c_x(p)=\frac12\sqrt{\phi_\infty(2x\log p)}\,\frac{1+q}{\sqrt q}. \tag{1.3}
\]
It satisfies
\[
 0<c_x(p)<1,\qquad
 c_x(p)^4=\frac{(1+q)^4}{8(1+q^4)}.                 \tag{1.4}
\]
The scalar cell has quadratic contact with its lossless boundary and therefore erases the
first-order logarithmic direction.  Its positive Julia defect
\[
 d_x(p)=\sqrt{1-c_x(p)^2}                            \tag{1.5}
\]
restores that direction.  The rotation
\[
 J_x(p)=
 \begin{pmatrix}
  c_x(p)&-d_x(p)\\
  d_x(p)& c_x(p)
 \end{pmatrix}                                      \tag{1.6}
\]
is the unique orientation-preserving real-plane isometry with the prescribed positive first
column, and one additional real channel is minimal.  Its normal derivative is arithmetic:
\[
 \ell_x(p):=\frac{2}{\sqrt3}\frac{d_x(p)}x
 \longrightarrow\log p\qquad(x\downarrow0).         \tag{1.7}
\]

The central point is that (1.2) and (1.7) are connected globally.  Define
\[
 L_{x,m}(d)=\sum_p d_p\ell_x(p)p^{-m}k_p.            \tag{1.8}
\]
On every finite packet, \(L_{x,m}(d)\to L_m(d)\).  Hence every family of bounded Hardy
operators extending \(L_{x,m}\) along a sequence \(x\downarrow0\) has norm tending to infinity.
The Julia coordinate is not a bounded workaround for (1.2); it is its canonical singular
normal cone.

The same normal coordinate has a full Mellin orbit.  On the prime-power cell \(n=p^r\), set
\[
 s_u(p,r)=
 \left(\frac{e^{-u\log n}}{n\log p}\right)^{1/2},
 \qquad h_{x,u}(p,r)=\ell_x(p)s_u(p,r).              \tag{1.9}
\]
Then
\[
 h_{x,u}(p,r)^2\longrightarrow
 \frac{\log p}{n}e^{-u\log n}.
\]
Extending by zero away from prime powers gives
\[
 M_u(n)=\frac{\Lambda(n)}n e^{-u\log n}.             \tag{1.10}
\]
Two faces of this one orbit are decisive:
\[
 M_{-1}(n)=\Lambda(n),\qquad
 M_{-1/2}(n)=\frac{\Lambda(n)}{\sqrt n}.             \tag{1.11}
\]

The first face supplies the hypertranscendental obstruction.  Let
\[
 P(s)=\sum_p p^{-s},\qquad
 \cK_{\rm val}=\C(s)\langle P(rs):r\ge1\rangle_{\rm diff}. \tag{1.12}
\]
On \(\Re s>1\), the Dirichlet carrier of \(M_{-1}\) is
\[
 V(s)=\sum_{n\ge1}\frac{\Lambda(n)}{n^s}
     =-\frac{\zeta'(s)}{\zeta(s)}.                  \tag{1.13}
\]
We prove that \(V\) is hypertranscendental over \(\cK_{\rm val}\).  Thus even access to every
prime-zeta valuation layer and all of their derivatives does not give a finite
differential-algebraic explanation of the Julia boundary carrier.

The second face supplies the divisor.  A finite fixed-width Weil scalarization of
\(\Lambda(n)/\sqrt n\), followed by the canonical Gamma--polar completion, produces a sequence
\(g_L(n)\) without using any zero of \(\Xi\).  The Dirichlet series of any admissible tail of
this sequence continues to \(\Re s>1\); its poles are exactly the shifted strict-right Xi zeros,
and its normalized residues are their analytic multiplicities.  Therefore one source sequence
recovers the complete off-critical divisor.

Finally, the divisor determines a nonnegative infinitesimal Julia-defect density.  Every
strict-right upper Schwarz pair has total mass \(8\pi\).  The constant is independent of the
zero's height and distance from the critical line.  This is an index law, not a smallness
estimate: the total may be infinite, and the theorem does not exclude a single pair.

The result closes four questions without claiming RH:
\begin{enumerate}[label=\textup{(\roman*)}]
\item which completion recovers the prime logarithmic clock after bounded Hardy closure fails;
\item whether that completion is minimal and whether its global singularity is forced;
\item whether the recovered arithmetic carrier remains outside every finite
differential-algebraic assembly of prime-zeta valuation layers;
\item whether the same completion reaches a multiplicity-complete global Riemann divisor
observable.
\end{enumerate}

Classical dilation theory supplies the abstract Julia operator
\cite{Halmos1950,SzNagyFoias2010}; no novelty is claimed for (1.6) in isolation.
Classical differential algebra supplies broad hypertranscendence criteria for Dirichlet series
\cite{Ostrowski1920,ShapiroSparer1986,Kolchin1973,AmouKatsurada2001,Pong2016,Nagoshi2016}.
Classical explicit formulae supply the prime--zero dictionary.  The new theorem is the forced
coupling: one arithmetic normal coordinate simultaneously realizes the only minimal positive
escape from the dense-graph wall, the full von Mangoldt Mellin orbit, a complete
valuation-layer Kolchin obstruction, a multiplicity-complete divisor transform, and a
quantized defect index.

\section{The arithmetic Julia normal-cone theorem}\label{sec:main}

Let \(\cP_0\) denote the real vector space of finitely supported prime packets.  All closures
below use the Hilbert norm of \(\cH\oplus\cH\).

\begin{definition}[Admissible tail]
For \(L>0\), an integer \(N\) is \(L\)-admissible when
\[
 L\le\frac12\log n\qquad(n\ge N).
\]
\end{definition}

Write every nontrivial Xi zero as \(\rho=\frac12+w\), with analytic multiplicity retained.
Let \(\cZ_+\) be the multiset of centered zeros with \(\Re w>0\), and let
\(\cZ_+^\uparrow\) contain one upper representative of each Schwarz pair.  For \(L>0\), define
\[
 A_L(w)=
 \begin{cases}
  (1-e^{-Lw})/w,&w\ne0,\\
  L,&w=0.
 \end{cases}                                        \tag{2.1}
\]

\begin{theorem}[Arithmetic Julia normal-cone law]\label{thm:normal-cone}
The source \(C\), clock \(L_m\), Gamma-sewn contraction \(c_x(p)\), and Julia normal coordinate
\(\ell_x(p)\) satisfy the following unconditional statements.

\begin{enumerate}[label=\textup{(\Alph*)}]
\item \textbf{Maximal Hardy singularity.}  For every \(m\ge0\),
\[
 \overline{\Graph(C,L_m)}=\cH\oplus\cH.              \tag{2.2}
\]
For every finite packet \(d\),
\[
 L_{x,m}(d)\longrightarrow L_m(d)\qquad(x\downarrow0). \tag{2.3}
\]
If \(x_j\downarrow0\) and \(F_j:\cH\to\cH\) are bounded operators satisfying
\[
 F_jC(d)=L_{x_j,m}(d)\qquad(j\ge1,\ d\in\cP_0),
\]
then
\[
 \|F_j\|\longrightarrow+\infty.                     \tag{2.4}
\]

\item \textbf{Unique minimal positive lift.}  At every \(x>0\) and prime \(p\), the scalar
map \(v\mapsto c_x(p)v\) is not an isometry.  The matrix \(J_x(p)\) in (1.6) is an isometry.
If \((a,b)\) is a unit vector orthogonal to \((c_x(p),d_x(p))\) and
\[
 c_x(p)b-d_x(p)a>0,
\]
then
\[
 (a,b)=(-d_x(p),c_x(p)).                             \tag{2.5}
\]
Thus one additional real channel is necessary and, after positive orientation, unique.

\item \textbf{Mellin--Kolchin boundary.}  Coefficientwise on \(\N\),
\[
 h_{x,u}^2\longrightarrow M_u
 \quad\text{for every }u\in\R,                      \tag{2.6}
\]
where \(M_u\) is given by (1.10).  The von Mangoldt carrier \(V\) in (1.13) is
hypertranscendental over \(\cK_{\rm val}\):
\[
 V\text{ satisfies no nonzero algebraic differential equation over }
 \C(s)\langle P(rs):r\ge1\rangle_{\rm diff}.         \tag{2.7}
\]

\item \textbf{Divisor completeness.}  For every \(L>0\) and every \(L\)-admissible \(N\),
there are source-defined finite-aperture coefficients \(g_{x,L,N}(n)\) with limits
\(g_{L,N}(n)\).  The ordinary Dirichlet series
\[
 \cD_{L,N}(s)=\sum_{n\ge N}\frac{g_{L,N}(n)}{n^s}
\]
has a canonical meromorphic continuation \(\cD_{L,N}^{\#}\) to \(\Re s>1\).  Its poles are
exactly
\[
 s=1+w,\qquad w\in\cZ_+,                            \tag{2.8}
\]
all are simple, and
\[
 \Res_{s=1+w}\cD_{L,N}^{\#}(s)=2m(w)A_L(w)^2.        \tag{2.9}
\]
In particular,
\[
 m(w)=\frac{\Res_{s=1+w}\cD_{L,N}^{\#}(s)}
 {2A_L(w)^2},                                       \tag{2.10}
\]
and
\[
 {\rm RH}\quad\Longleftrightarrow\quad
 \cD_{L,N}^{\#}\text{ is holomorphic on }\Re s>1.   \tag{2.11}
\]

\item \textbf{Quantized divisor index.}  There is a nonnegative infinitesimal defect density
\(\mathfrak d(t)\) on the Euler edge such that
\[
 \int_{\R}\mathfrak d(t)\,dt
 =\sum_{q\in\cZ_+^\uparrow}8\pi,                    \tag{2.12}
\]
where the equality is understood in \([0,+\infty]\) and multiplicity is retained.  Every
summand is independent of the pair's height and transverse displacement.  Consequently
\[
 {\rm RH}\quad\Longleftrightarrow\quad
 \int_{\R}\mathfrak d(t)\,dt=0.                     \tag{2.13}
\]
\end{enumerate}
\end{theorem}

\begin{remark}[Logical content]
Equations (2.11) and (2.13) are exact classifiers, not proofs that the right sides hold.  The
unconditional advance is the construction and rigidity of the source-to-divisor channel:
bounded Hardy geometry cannot realize it, finite differential algebra over all valuation layers
cannot synthesize its Euler face, yet its critical face is lossless for the complete off-critical
divisor.
\end{remark}

The proof occupies Sections \ref{sec:hardy}--\ref{sec:index}.

\section{The bounded-completion wall}\label{sec:hardy}

For \(f=(f_0,f_1,\ldots)\in\cH\), write
\[
 F_f(z)=\sum_{j\ge0}f_jz^j,\qquad |z|<1.
\]
Then \(\langle k_p,f\rangle=F_f(p^{-1})\).

\begin{lemma}[Prime logarithmic germ rigidity]\label{lem:germ-rigidity}
Fix \(m\ge0\).  If \(f,g\in\cH\) satisfy
\[
 F_f(p^{-1})+(\log p)p^{-m}F_g(p^{-1})=0           \tag{3.1}
\]
for every prime \(p\), then \(f=g=0\).
\end{lemma}

\begin{proof}
Let \(Bf=(f_1,f_2,\ldots)\) be the backward shift.  The identity
\[
 F_f(z)=f_0+zF_{Bf}(z)                              \tag{3.2}
\]
holds on the unit disc.  Since \(p^{-1}\to0\), \(\log p\to+\infty\), and
\((\log p)p^{-1}\to0\), equation (3.1) with \(m\ge1\) first gives \(f_0=0\).
Dividing by \(p^{-1}\) and using (3.2) reduces order \(m\) to \(m-1\).  Iteration removes the
first \(m\) coefficients of \(f\) and reduces to
\[
 F_u(p^{-1})+(\log p)F_g(p^{-1})=0.                 \tag{3.3}
\]

Divide (3.3) by \(\log p\) and pass through the primes to obtain \(g_0=0\).  Substituting
\(F_g(z)=zF_{Bg}(z)\) back into (3.3) and using
\((\log p)p^{-1}\to0\) gives \(u_0=0\).  Division by \(p^{-1}\) now shows that
\((Bu,Bg)\) satisfies the same order-zero relation.  Induction gives every Taylor coefficient
of \(u\) and \(g\) equal to zero.  The previously removed coefficients of \(f\) were already
zero, hence \(f=g=0\).
\end{proof}

\begin{proposition}[Dense graph]\label{prop:dense-graph}
For every \(m\ge0\), equation \emph{(2.2)} holds.
\end{proposition}

\begin{proof}
An element \((f,g)\in\cH\oplus\cH\) orthogonal to the graph satisfies
\[
 \langle k_p,f\rangle+
 (\log p)p^{-m}\langle k_p,g\rangle=0
\]
for every prime \(p\), by testing the singleton packet at \(p\).  Lemma
\ref{lem:germ-rigidity} gives \(f=g=0\).  Thus the graph has trivial orthogonal complement,
and its closure is the full product.
\end{proof}

The next consequence makes the singularity quantitative.

\begin{proposition}[Bounded recovery blow-up]\label{prop:blowup}
Let \(K\) be a real Hilbert space and \(B:\cH\to K\) a nonzero bounded operator.  If bounded
operators \(F_j:\cH\to K\) satisfy
\[
 F_jC(d)\longrightarrow BL_m(d)
\]
for every \(d\in\cP_0\), then \(\|F_j\|\to+\infty\).
\end{proposition}

\begin{proof}
First observe that no finite \(M\) can satisfy
\[
 \|BL_m(d)\|\le M\|C(d)\|\qquad(d\in\cP_0).         \tag{3.4}
\]
Indeed, Proposition \ref{prop:dense-graph} implies that the closure of the postprocessed graph
contains \((0,By)\) for every \(y\in\cH\).  Choose \(y\) with \(By\ne0\).  The closed cone
\(\{(u,v):\|v\|\le M\|u\|\}\) cannot contain \((0,By)\), contradicting (3.4).

Given a threshold \(b\), choose \(M>b\) and a packet \(d\) violating (3.4).  Packetwise
convergence implies eventually
\[
 M\|C(d)\|<\|F_jC(d)\|\le\|F_j\|\,\|C(d)\|.
\]
The violating packet has \(C(d)\ne0\), so \(\|F_j\|>M>b\) eventually.  Since \(b\) was
arbitrary, \(\|F_j\|\to+\infty\).
\end{proof}

\section{The Gamma-sewn Julia normal coordinate}\label{sec:julia}

Let \(x>0\), \(p\) be prime, and \(q=e^{-x\log p}\).  Define \(c_x(p)\) by (1.3) and
\(d_x(p)\) by (1.5).

\begin{proposition}[Exact contraction law]\label{prop:contraction}
For every \(x>0\) and prime \(p\),
\[
 0<c_x(p)<1,\qquad c_x(p)^2+d_x(p)^2=1,             \tag{4.1}
\]
and
\[
 c_x(p)^4=\frac{(1+q)^4}{8(1+q^4)}.                \tag{4.2}
\]
Moreover,
\[
 d_x(p)=\frac{\sqrt3}{2}x\log p+o(x)
 \qquad(x\downarrow0).                              \tag{4.3}
\]
\end{proposition}

\begin{proof}
The identities
\[
 \phi_\infty(2x\log p)^2=\frac{2q^2}{1+q^4}
\]
and (1.3) give (4.2).  Positivity is immediate.  Since
\[
 8(1+q^4)-(1+q)^4=(1-q)^2(7q^2+10q+7),             \tag{4.4}
\]
one has \(c_x(p)<1\) when \(0<q<1\).  Equation (4.1) follows from the definition of \(d_x(p)\).
Finally \(q=1-x\log p+O(x^2)\); expanding (4.2) at \(q=1\) gives
\[
 1-c_x(p)^2=\frac34x^2(\log p)^2+o(x^2).
\]
Taking the positive square root proves (4.3).
\end{proof}

\begin{proposition}[Unique minimal positive Julia lift]\label{prop:minimal}
The assertions in Theorem \ref{thm:normal-cone}\textup{(B)} hold.
\end{proposition}

\begin{proof}
Because \(0<c_x(p)<1\), the scalar map cannot preserve the norm of \(1\); hence no
one-dimensional isometric completion retaining that scalar coordinate exists.  Equation (4.1)
shows directly that \(J_x(p)\) is orthogonal.

For uniqueness, let \(c=c_x(p)\), \(d=d_x(p)\).  If \((a,b)\) is a unit vector orthogonal to
\((c,d)\), then
\[
 (cb-da)^2=(c^2+d^2)(a^2+b^2)-(ca+db)^2=1.
\]
Positive orientation gives \(cb-da=1\).  Therefore
\[
 a=-d(cb-da)+c(ca+db)=-d,\qquad
 b=c(cb-da)+d(ca+db)=c.
\]
\end{proof}

\begin{corollary}[Literal clock normal derivative]\label{cor:clock}
For every prime \(p\),
\[
 \ell_x(p)=\frac{2}{\sqrt3}\frac{d_x(p)}x
 \longrightarrow\log p.
\]
\end{corollary}

\begin{proposition}[Global packet bridge]\label{prop:packet}
For every \(m\ge0\) and \(d\in\cP_0\),
\[
 L_{x,m}(d)\longrightarrow L_m(d)
\quad\text{in }\cH.
\]
If \(x_j\downarrow0\) and bounded \(F_j:\cH\to\cH\) satisfy
\[
 F_jC(d)=L_{x_j,m}(d)
\]
for all \(j\) and all finite packets, then \(\|F_j\|\to+\infty\).
\end{proposition}

\begin{proof}
The first assertion is a finite sum of the scalar limits in Corollary \ref{cor:clock}.  The
second follows by composing those limits with \(x_j\downarrow0\) and applying Proposition
\ref{prop:blowup} with \(B\) equal to the identity.
\end{proof}

\begin{remark}[Rank]
Minimality is fiberwise: one additional real channel is necessary on each prime cell.  The
global prime-indexed completion is not claimed to be a single rank-one operator on \(\cH\).
Proposition \ref{prop:packet} proves instead the exact global statement available in the Hardy
topology: finite packets converge, while every bounded global realization escapes to infinite
operator norm.
\end{remark}

\section{One Mellin orbit and the valuation-layer Kolchin wall}\label{sec:kolchin}

For \(n=p^r\), \(r\ge1\), define \(s_u(p,r)\) and \(h_{x,u}(p,r)\) by (1.9), and set the squared
coefficient to zero when \(n\) is not a prime power.

\begin{proposition}[Julia--Mellin boundary law]\label{prop:mellin}
For every \(u\in\R\) and \(n\ge1\),
\[
 h_{x,u}(n)^2\longrightarrow
 M_u(n)=\frac{\Lambda(n)}n e^{-u\log n}.             \tag{5.1}
\]
The orbit obeys
\[
 M_{u+v}(n)=e^{-v\log n}M_u(n),                     \tag{5.2}
\]
with distinguished faces
\[
 M_{-1}=\Lambda,\qquad M_{-1/2}(n)=\Lambda(n)n^{-1/2}. \tag{5.3}
\]
\end{proposition}

\begin{proof}
If \(n=p^r\), Corollary \ref{cor:clock} and (1.9) give
\[
 h_{x,u}(p,r)^2\longrightarrow
 (\log p)^2\frac{e^{-u\log n}}{n\log p}
 =\frac{\Lambda(n)}n e^{-u\log n}.
\]
Both sides vanish away from prime powers.  Equations (5.2)--(5.3) are immediate.
\end{proof}

We now show that the Euler face is not a finite differential-algebraic recombination of all
prime-zeta valuation layers.  The statement is relative, and is substantially stronger than
hypertranscendence over \(\C(s)\).

Consider the vector space of complex prime-power profiles modulo profiles supported on finitely
many prime rays.  The profile of a prime-power-supported arithmetic source \(a\) is
\[
 \widehat a(p,r)=a(p^r).
\]
Differentiation of the corresponding Dirichlet carrier acts on profiles by the clock
\[
 (T\widehat a)(p,r)=r\log p\,\widehat a(p,r).        \tag{5.4}
\]
The \(r\)-th valuation layer has profile
\[
 e_r(p,k)=\mathbf 1_{\{k=r\}}
\]
and analytic carrier \(P(rs)\).

\begin{lemma}[Unused-layer separation]\label{lem:unused}
Let \(a\) be supported on positive prime powers and suppose that for infinitely many \(r\),
\[
 a(p^r)\ne0\qquad\text{for every prime }p.           \tag{5.5}
\]
Then the complete clock orbit of \(a\) is linearly independent, modulo finite prime-ray support,
from the clock orbits of every finite family of valuation layers.
\end{lemma}

\begin{proof}
Fix layers \(r_1,\ldots,r_k\), a clock order \(J\), and a linear relation involving
\(T^ja\) and \(T^je_{r_i}\), \(0\le j\le J\), modulo finitely many prime rays.  Choose a layer
\(r_0\) satisfying (5.5) and distinct from every \(r_i\).  Restriction to \(r_0\) kills all old
layers and leaves
\[
 a(p^{r_0})\sum_{j=0}^J c_j(r_0\log p)^j=0
\]
for all but finitely many primes.  The first factor never vanishes.  Since \(\log p\) takes
infinitely many distinct values, the polynomial is zero and every \(c_j=0\).  Returning to
layer \(r_i\) and using the same polynomial argument kills the coefficients of each old layer.
\end{proof}

\begin{lemma}[Carrier transfer]\label{lem:carrier}
For finite-abscissa prime-power sources, clock-orbit independence modulo finitely many prime rays
implies differential independence of their Dirichlet carriers over \(\C(s)\).
\end{lemma}

\begin{proof}
This is the prime-profile specialization of the Jacobian criterion for Dirichlet series
\cite{ShapiroSparer1986}.  We recall the mechanism.  Any alleged algebraic relation among
finitely many jets can be differentiated minimally until its Jacobian in a highest jet is
nonzero.  Dirichlet-series uniqueness transfers that Jacobian to the convolution algebra.
Prime-ray derivations then isolate its least active prime-power profiles.  A nonzero Jacobian
would give a finite linear relation among the corresponding clock jets modulo finitely many
prime rays, contrary to the hypothesis.  Adjoining the genuine coordinate \(s\) creates no new
relation: after clearing denominators, a polynomial relation
\[
 \sum_{j=0}^d s^j\sum_{n\ge1}b_{j,n}n^{-s}=0
\]
is killed by selecting the least active \(n\) and sending real \(s\to+\infty\).
\end{proof}

\begin{theorem}[Complete valuation-layer hypertranscendence]\label{thm:ht}
The von Mangoldt carrier
\[
 V(s)=\sum_{n\ge1}\Lambda(n)n^{-s}
\]
is hypertranscendental over
\[
 \cK_{\rm val}=\C(s)\langle P(rs):r\ge1\rangle_{\rm diff}.
\]
\end{theorem}

\begin{proof}
The von Mangoldt source is supported on positive prime powers and
\[
 \Lambda(p^r)=\log p\ne0
\]
for every prime \(p\) and every \(r\ge1\).  Lemma \ref{lem:unused} therefore applies after any
finite choice of valuation layers.  Lemma \ref{lem:carrier} makes \(V\) differentially
independent over the differential field generated by those finitely many carriers.

Every differential polynomial over \(\cK_{\rm val}\) uses only finitely many elements of that
field; after writing their numerators and denominators, only finitely many layer generators and
jets occur.  The preceding finite statement excludes the relation.  Hence no nonzero
differential polynomial over the full field annihilates \(V\).  Absolute convergence on
\(\Re s>1\) gives \(V=-\zeta'/\zeta\).
\end{proof}

\begin{remark}[What is and is not proved]
Theorem \ref{thm:ht} concerns the Euler face \(u=-1\).  The critical face \(u=-1/2\) belongs to
the same source-defined Mellin orbit, but we do not claim that it is hypertranscendental over the
unshifted field \(\cK_{\rm val}\).  Its role is instead the exact divisor transgression of the
next section.
\end{remark}

\section{The critical face reconstructs the divisor}\label{sec:divisor}

Fix \(L>0\) and \(a\ge L\).  Let
\[
 h_{L,a}(t)=\mathbf1_{[a-L,a]}(t)+
             \mathbf1_{[-a,-a+L]}(t)                \tag{6.1}
\]
and let
\[
 R_{L,a}(t)=\int_{\R}h_{L,a}(y+t)h_{L,a}(y)\,dy.    \tag{6.2}
\]
The autocorrelation is continuous, piecewise linear, and supported on \([-2a,2a]\).  Define
\[
 F_{a,L}(w)=e^{aw}A_L(w)+e^{-aw}A_L(-w)             \tag{6.3}
\]
and the absolutely convergent Xi polarization
\[
 Q_L(a)=\sum_{\Xi(1/2+w)=0}F_{a,L}(w)^2,            \tag{6.4}
\]
with analytic multiplicity.

At finite aperture put
\[
 \cP_x(L,a)=
 \sum_{2\le n\le e^{2a}}
 h_{x,-1/2}(n)^2R_{L,a}(\log n).                    \tag{6.5}
\]
The sum is finite.  Separate the completed factor
\[
 \cG(s)=\frac12s(s-1)\pi^{-s/2}\Gamma(s/2)
\]
and define
\[
 \cA_{\infty,L,a}=
 \frac1\pi\int_{\R}
 \widehat h_{L,a}(t-i)^2
 \frac{\cG'}{\cG}\!\left(\frac32+it\right)\,dt,     \tag{6.6}
\]
where \(\widehat h_{L,a}(z)=F_{a,L}(iz)\).  Finally set
\[
 g_{x,L}(a)=-2\cP_x(L,a)+\Re\cA_{\infty,L,a}.       \tag{6.7}
\]
Every quantity in (6.5)--(6.7) is constructed from primes, the Gamma factor, and the fixed
window; no Xi zero enters its definition.

\begin{proposition}[Julia-to-Xi transgression]\label{prop:transgression}
For \(0<L\le a\),
\[
 g_{x,L}(a)\longrightarrow Q_L(a)
\qquad(x\downarrow0).                               \tag{6.8}
\]
\end{proposition}

\begin{proof}
Proposition \ref{prop:mellin} at \(u=-1/2\) and finiteness of (6.5) give
\[
 \cP_x(L,a)\longrightarrow
 \sum_{2\le n\le e^{2a}}
 \frac{\Lambda(n)}{\sqrt n}R_{L,a}(\log n).         \tag{6.9}
\]
The sharp Guinand--Weil formula for (6.1) is
\[
 Q_L(a)=
 -2\sum_{2\le n\le e^{2a}}
 \frac{\Lambda(n)}{\sqrt n}R_{L,a}(\log n)
 +\Re\cA_{\infty,L,a}.                              \tag{6.10}
\]
For completeness, (6.10) follows by integrating
\[
 \widehat h_{L,a}(z)^2
 \frac{\Xi_C'(z)}{\Xi_C(z)},\qquad
 \Xi_C(z)=\Xi\!\left(\frac12+iz\right),
\]
over rectangles with horizontal edges at imaginary parts \(\pm1\).
The residue sum is (6.4).  The vertical edges vanish by the centered Hadamard product and the
\(O(|z|^{-1})\) decay of \(\widehat h_{L,a}\).  On the remaining edge,
\[
 \frac{\zeta'}{\zeta}\!\left(\frac32+it\right)
 =-\sum_{n\ge2}\frac{\Lambda(n)}{n^{3/2}}e^{-it\log n}.
\]
Fourier inversion yields
\[
 \frac1\pi\int_{\R}\widehat h_{L,a}(t-i)^2e^{-ity}\,dt
 =2e^yR_{L,a}(y).
\]
Because \(R_{L,a}\) has support in \([-2a,2a]\), the prime sum is exactly finite.  This proves
(6.10), and (6.8) follows from (6.9).
\end{proof}

Sample at \(a=\frac12\log n\).  Choose an \(L\)-admissible \(N\), and set
\[
 g_{x,L,N}(n)=
 \begin{cases}
  g_{x,L}(\frac12\log n),&n\ge N,\\
  0,&n<N,
 \end{cases}
\qquad
 g_{L,N}(n)=
 \begin{cases}
  Q_L(\frac12\log n),&n\ge N,\\
  0,&n<N.
 \end{cases}                                       \tag{6.11}
\]
Proposition \ref{prop:transgression} gives coefficientwise convergence.

\begin{proposition}[Multiplicity-complete continuation]\label{prop:tomography}
The continuation and residue assertions \emph{(2.8)--(2.11)} hold for the sequence in
\emph{(6.11)}.
\end{proposition}

\begin{proof}
Expand (6.3) and fold the centered divisor by \(w\mapsto-w\).  The moving right edge of
\(Q_L(\frac12\log n)\) is
\[
 2\sum_{\substack{w\in\cZ\\ \Re w>0}}
 m(w)A_L(w)^2n^w,                                  \tag{6.12}
\]
while the remaining self, left-edge, and center terms form a bounded coefficient sequence.
Initially in a right half-plane, the Dirichlet series of (6.12) is
\[
 2\sum_{\substack{w\in\cZ\\ \Re w>0}}
 m(w)A_L(w)^2\zeta(s-w),                            \tag{6.13}
\]
after subtraction of an entire finite Dirichlet polynomial coming from \(n<N\).
Normal convergence away from the candidate poles continues (6.13) to \(\Re s>1\); the bounded
remainder is holomorphic there.

Each \(w\) produces a simple pole at \(s=1+w\) with residue \(2m(w)A_L(w)^2\).  There is no
cancellation: if several zeros have the same \(w\), they have already been combined into the
positive analytic multiplicity \(m(w)\), and \(A_L(w)\ne0\) for \(\Re w>0\).
This proves (2.8)--(2.10).  Functional-equation and Schwarz symmetry reconstruct every
off-critical zero from the strict-right divisor.  Therefore the continuation is holomorphic
throughout \(\Re s>1\) exactly when no strict-right zero exists, which is equivalent to RH.
\end{proof}

\begin{remark}[Why this is not a decorative equivalence]
The sequence in (6.11) is not defined from the target divisor.  Its finite-aperture precursor is
the squared Julia environment coordinate in (6.5), and every coefficient uses only finitely many
von Mangoldt weights plus the fixed Gamma--polar integral.  Proposition
\ref{prop:transgression} is the source-to-divisor map; Proposition \ref{prop:tomography} proves
that this map loses neither location nor multiplicity.
\end{remark}

\section{The \texorpdfstring{\(8\pi\)}{8pi} defect index}\label{sec:index}

For a strict-right upper centered zero \(\lambda=\beta+i\gamma\), \(\beta>0\), the associated
continued Ward pair response on the Euler edge has real part
\[
 r_\lambda(t)=
 \frac{2(\beta+\frac12)}
 {(\beta+\frac12)^2+(t+\gamma)^2}
 +
 \frac{2(\beta+\frac12)}
 {(\beta+\frac12)^2+(t-\gamma)^2}.                  \tag{7.1}
\]
Both terms are nonnegative Lorentzians.  Let \(r(t)\) be the sum over all upper Schwarz pairs,
with multiplicity, and define the absolutely convergent complex response on \(\Re w>0\) by
\[
 \mathcal R(w)=2\sum_{\lambda\in\cZ_+^\uparrow}
 \left((w+\lambda)^{-1}+(w+\overline\lambda)^{-1}\right). \tag{7.2}
\]
Thus \(r(t)=\Re\mathcal R(\frac12+it)\).

This response is the first variation of an explicit contractive phase.  Let \(E_+\) be the
genus-one canonical product over the strict-right centered divisor, paired by Schwarz symmetry,
and put
\[
 B=2\sum_{\lambda\in\cZ_+^\uparrow}
 \left(\lambda^{-1}+\overline\lambda^{-1}\right).   \tag{7.3}
\]
Schwarz pairing makes this series absolutely convergent.  For \(\Re w>0\) and real aperture
\(a\ge0\), define
\[
 \Phi_w(a)=
 \frac{E_+(iw-ai)}{E_+(iw+ai)}e^{-aB}.              \tag{7.4}
\]
Cancellation of the paired genus-one contacts writes \(\Phi_w(a)\) as the absolutely
convergent product of the Cayley factors
\[
 \frac{w+\lambda-a}{w+\lambda+a}
 \frac{w+\overline\lambda-a}{w+\overline\lambda+a}.
\]
Every factor has norm at most one, so \(|\Phi_w(a)|\le1\).  If
\[
 \Delta_w(a)=1-|\Phi_w(a)|^2,
\]
then \(\Phi_w(0)=1\) and
\[
 \left.\frac{d}{da}\Delta_w(a)\right|_{a=0}
 =2\Re\mathcal R(w).                                 \tag{7.5}
\]
Define the Euler-edge infinitesimal defect density by
\[
 \mathfrak d(t)=2r(t).                              \tag{7.6}
\]

\begin{proposition}[Universal pair charge]\label{prop:charge}
Every strict-right upper Schwarz pair has
\[
 \int_{\R}2r_\lambda(t)\,dt=8\pi.                   \tag{7.7}
\]
Consequently (2.12)--(2.13) hold.
\end{proposition}

\begin{proof}
For \(a>0\),
\[
 \int_{\R}\frac{2a}{a^2+(t-c)^2}\,dt=2\pi.          \tag{7.8}
\]
Equation (7.1) is the sum of two such profiles with
\(a=\beta+\frac12>0\), so its mass is \(4\pi\), independent of \(\beta\) and \(\gamma\).
The defect derivative in (7.5) contributes the additional factor \(2\), proving (7.7).

Nonnegativity permits Tonelli's theorem:
\[
 \int_{\R}\mathfrak d(t)\,dt
 =\sum_{\lambda\in\cZ_+^\uparrow}
 \int_{\R}2r_\lambda(t)\,dt
 =\sum_{\lambda\in\cZ_+^\uparrow}8\pi.
\]
The integral vanishes exactly when the indexing multiset is empty.  Functional-equation and
Schwarz symmetry identify emptiness of the strict-right multiset with RH.
\end{proof}

\begin{remark}[No zero-exclusion inference]
The charge is quantized but not bounded above by a finite source budget in this theorem.
Neither high zeros nor zeros close to the critical line have smaller total charge; they only
redistribute the Lorentzian profile.  Equation (7.7) therefore supplies a robust divisor index,
not a contradiction from one hypothetical zero.
\end{remark}

\begin{proof}[Proof of Theorem \ref{thm:normal-cone}]
Part (A) is Propositions \ref{prop:dense-graph} and \ref{prop:packet}.  Part (B) is Proposition
\ref{prop:minimal}.  Part (C) is Proposition \ref{prop:mellin} and Theorem \ref{thm:ht}.
Part (D) is Propositions \ref{prop:transgression} and \ref{prop:tomography}.  Part (E) is
Proposition \ref{prop:charge}.
\end{proof}

\section{Relation to existing RH routes}\label{sec:literature}

The theorem intersects four established subjects, but its novelty is not the union of four
citations.

\paragraph{Dilation theory.}
The Julia operator and positive defect operator are classical
\cite{Halmos1950,SzNagyFoias2010}.  The new assertion is arithmetic and singular: the Gamma-sewn
compression selects a normal derivative equal to \(\log p\), finite packets converge to the
actual-prime logarithmic clock, and every bounded global Hardy realization must diverge in
operator norm.  Abstract dilation theory alone contains none of the prime-clock, Mellin, or
divisor conclusions.

\paragraph{Differential algebra.}
Ostrowski initiated differential transcendence for Dirichlet series; Shapiro--Sparer and later
work developed algebraic and differential independence criteria
\cite{Ostrowski1920,ShapiroSparer1986,AmouKatsurada2001,Pong2016,Nagoshi2016}.
Theorem \ref{thm:ht} is relative to the single field containing the genuine coordinate and every
prime-zeta valuation layer at once.  Its role is not to restate the classical
hypertranscendence of zeta.  It proves that the precise Euler boundary generated by the Julia
normal coordinate cannot be produced by any finite differential-algebraic circuit over the
complete layer field.

\paragraph{Explicit formula and localization.}
The Guinand--Weil--Barner formula, Weil positivity, and localized Weil forms are classical or
actively developed \cite{Guinand1949,Weil1952,Barner1981,Bombieri2000,Bombieri2003,Suzuki2026}.
We claim no new general explicit formula.  New content is upstream: the critical coefficient
\(\Lambda(n)/\sqrt n\) is obtained as a boundary face of the unique minimal Julia environment,
and the finite scalarization of that environment is shown to retain the entire off-critical
divisor with multiplicity.

\paragraph{Spectral and semilocal programs.}
Connes and collaborators construct trace-formula and zeta-cycle structures, and recent zeta
spectral triples use finite-prime self-adjoint rank-one perturbations whose rigorous limiting
identification would imply RH
\cite{Connes1999,ConnesConsani2021,ConnesConsani2023,CCM2025}.
Our result neither proves nor disproves that convergence.  It supplies a complementary
source-side law: any Hardy-pivot realization faithful to the actual-prime logarithmic clock must
be singular, while the canonical singular normal coordinate already reaches a
multiplicity-complete divisor observable.  Here ``one extra channel'' is fiberwise at each
prime, not a claim that the global construction is a rank-one perturbation of a CCM operator.

Searches across these literatures locate the ingredients separately, but not the following
causal chain:
\[
\boxed{
\begin{gathered}
\text{Gamma-sewn prime contraction}
\longrightarrow
\text{unique positive Julia normal clock}\\
\longrightarrow
\text{forced global Hardy singularity}
\longrightarrow
\text{one von Mangoldt Mellin orbit}\\
\longrightarrow
\begin{cases}
\text{complete valuation-layer hypertranscendence},\\
\text{multiplicity-complete Riemann divisor tomography}
\end{cases}\\
\longrightarrow
\text{\(8\pi\) defect index per strict-right pair}.
\end{gathered}}
\]

\section{Consequences and closed scope}\label{sec:scope}

The theorem gives a universal completion law with five sharp boundaries.

\begin{enumerate}[label=\textup{(\arabic*)}]
\item \textbf{Topological boundary.}  Finite trace subtraction and arbitrary bounded
postprocessing do not make the prime logarithmic clock closable.  Faithful bounded
regularizations must leave every operator-norm ball.
\item \textbf{Dimensional boundary.}  On each source cell, the positive Julia environment is
the unique minimal isometric completion.  Its first-order normal coordinate is exactly the
prime logarithm after one universal normalization.
\item \textbf{Differential-algebraic boundary.}  The Euler face remains
hypertranscendental even after every prime-zeta valuation layer is adjoined.
\item \textbf{Information boundary.}  The critical face, after canonical Gamma--polar
completion, reconstructs every off-critical zero and analytic multiplicity.  No source
information relevant to the horizontal divisor is lost.
\item \textbf{Index boundary.}  The reconstructed strict-right divisor is measured by a
nonnegative index with universal quantum \(8\pi\).
\end{enumerate}

This opens a proof route but does not finish it.  What remains is not an unidentified open
question inside the theorem.  It is the distinct RH task of proving that the explicitly
constructed divisor-complete continuation is holomorphic on \(\Re s>1\), equivalently that its
nonnegative defect index vanishes.  The present work answers the completion problem completely:
it identifies the minimal source channel, proves the singularity any realization must possess,
shows why finite differential algebra cannot replace it, and proves that its critical boundary
contains the full target divisor.

The law is universal in three senses.  It holds at every finite Hardy trace order, for every
positive aperture and every prime cell, and at every positive fixed window width and admissible
tail.  It uses no rightmost zero, simplicity assumption, numerical zero verification, finite
off-line-zero hypothesis, or positivity conjecture.

\begin{thebibliography}{99}
\small

\bibitem{AmouKatsurada2001}
M.~Amou and M.~Katsurada,
\newblock Differential transcendence of a class of generalized Dirichlet series,
\newblock \emph{Illinois J. Math.} \textbf{45} (2001), 939--948,
\newblock \href{https://doi.org/10.1215/ijm/1258138161}{doi:10.1215/ijm/1258138161}.

\bibitem{Barner1981}
K.~Barner,
\newblock On A. Weil's explicit formula,
\newblock \emph{J. Reine Angew. Math.} \textbf{323} (1981), 139--152.

\bibitem{Bombieri2000}
E.~Bombieri,
\newblock Remarks on Weil's quadratic functional in the theory of prime numbers, I,
\newblock \emph{Rend. Lincei Mat. Appl.} \textbf{11} (2000), 183--233.

\bibitem{Bombieri2003}
E.~Bombieri,
\newblock A variational approach to the explicit formula,
\newblock \emph{Comm. Pure Appl. Math.} \textbf{56} (2003), 1151--1164.

\bibitem{CCM2025}
A.~Connes, C.~Consani, and H.~Moscovici,
\newblock Zeta spectral triples,
\newblock arXiv:2511.22755, 2025.

\bibitem{Connes1999}
A.~Connes,
\newblock Trace formula in noncommutative geometry and the zeros of the Riemann zeta function,
\newblock \emph{Selecta Math. (N.S.)} \textbf{5} (1999), 29--106.

\bibitem{ConnesConsani2021}
A.~Connes and C.~Consani,
\newblock Weil positivity and trace formula, the archimedean place,
\newblock \emph{Selecta Math. (N.S.)} \textbf{27} (2021), article 77.

\bibitem{ConnesConsani2023}
A.~Connes and C.~Consani,
\newblock Spectral triples and zeta-cycles,
\newblock \emph{Enseign. Math.} \textbf{69} (2023), 93--148.

\bibitem{Guinand1949}
A.~P. Guinand,
\newblock Fourier reciprocities and the Riemann zeta-function,
\newblock \emph{Proc. London Math. Soc. (2)} \textbf{51} (1949), 401--414.

\bibitem{Halmos1950}
P.~R. Halmos,
\newblock Normal dilations and extensions of operators,
\newblock \emph{Summa Brasil. Math.} \textbf{2} (1950), 125--134.

\bibitem{Kolchin1973}
E.~R. Kolchin,
\newblock \emph{Differential Algebra and Algebraic Groups},
\newblock Academic Press, New York, 1973.

\bibitem{Nagoshi2016}
H.~Nagoshi,
\newblock Hypertranscendence of \(L\)-functions for
\(\mathrm{GL}_m(\mathbb A_{\mathbb Q})\),
\newblock \emph{Bull. Aust. Math. Soc.} \textbf{93} (2016), 388--399,
\newblock \href{https://doi.org/10.1017/S000497271500129X}
{doi:10.1017/S000497271500129X}.

\bibitem{Ostrowski1920}
A.~Ostrowski,
\newblock \"Uber Dirichletsche Reihen und algebraische Differentialgleichungen,
\newblock \emph{Math. Z.} \textbf{8} (1920), 241--298,
\newblock \href{https://doi.org/10.1007/BF01206523}{doi:10.1007/BF01206523}.

\bibitem{Pong2016}
W.~Y. Pong,
\newblock Applications of differential algebra to algebraic independence of arithmetic
functions,
\newblock \emph{Acta Arith.} \textbf{172} (2016), 149--173,
\newblock \href{https://doi.org/10.4064/aa8112-12-2015}{doi:10.4064/aa8112-12-2015}.

\bibitem{ShapiroSparer1986}
H.~N. Shapiro and G.~H. Sparer,
\newblock On algebraic independence of Dirichlet series,
\newblock \emph{Comm. Pure Appl. Math.} \textbf{39} (1986), 695--745,
\newblock \href{https://doi.org/10.1002/cpa.3160390602}
{doi:10.1002/cpa.3160390602}.

\bibitem{Suzuki2025}
M.~Suzuki,
\newblock On the Hilbert space derived from the Weil distribution,
\newblock \emph{Canad. J. Math.}, published online 3 November 2025,
\newblock \href{https://doi.org/10.4153/S0008414X25101739}
{doi:10.4153/S0008414X25101739}.

\bibitem{Suzuki2026}
M.~Suzuki,
\newblock Weil's quadratic form via the screw function,
\newblock arXiv:2606.09096, 2026.

\bibitem{SzNagyFoias2010}
B.~Sz.-Nagy, C.~Foias, H.~Bercovici, and L.~K\'erchy,
\newblock \emph{Harmonic Analysis of Operators on Hilbert Space}, 2nd ed.,
\newblock Springer, New York, 2010,
\newblock \href{https://doi.org/10.1007/978-1-4419-6094-8}
{doi:10.1007/978-1-4419-6094-8}.

\bibitem{Weil1952}
A.~Weil,
\newblock Sur les ``formules explicites'' de la th\'eorie des nombres premiers,
\newblock \emph{Commun. S\'em. Math. Univ. Lund}, volume dedicated to Marcel Riesz
(1952), 252--265.

\end{thebibliography}

\end{document}
